\chapter{Literature Review}
\label{chap:litreview}

This chapter reviews the theoretical foundations underlying the study and
situates the contribution within the existing literature.  Seven areas
are covered: single-echelon inventory theory, Clark-Scarf multi-echelon
optimality, periodic-review and cyclic policies, the bullwhip effect,
transport consolidation, simulation-optimisation methods, and existing
open-source MEIO software.

%% -----------------------------------------------------------------------
\section{Single-Echelon Foundations}
\label{sec:lit:singleechelon}

\subsection{Newsvendor model}

The \textbf{newsvendor model} \citep{ArrowHarrisMarschak1951} is the
canonical single-period stocking problem.  A retailer must commit to a
stock level $S$ before random demand $D$ is realised; unsold units cost
$h$ per unit (holding) while unmet demand costs $p$ per unit (shortage).
The optimal stock level equates the critical ratio:

\begin{equation}
  F(S^*) = \frac{p}{p + h}
  \label{eq:newsvendor-crit}
\end{equation}

where $F$ is the CDF of demand.  For Gaussian demand this yields the
service-level formula $S^* = \mu + z_\alpha \sigma$, where $z_\alpha$ is
the quantile of the standard normal corresponding to the target fill rate
$\alpha$.

\subsection{Base-stock policy}

The \textbf{base-stock (order-up-to) policy} extends the newsvendor to a
multi-period setting with backlogging.  Each period the stocking point
orders up to a fixed target $S$, which is also the \emph{inventory
position} target.  Under i.i.d. demand and deterministic lead time $L$
\citep{Zipkin2000}, the optimal base-stock level is:

\begin{equation}
  S^* = \mu(L+1) + z_\alpha \sigma\sqrt{L+1}
  \label{eq:bs-optimal}
\end{equation}

This formula is used to initialise the analytical baseline in
experiment E0 (Section~\ref{sec:results:e0}).

\subsection{$(s,S)$ and $(R,Q)$ policies}

The \textbf{$(s,S)$ policy} places an order to restore inventory to level
$S$ whenever the inventory position falls below $s$
\citep{ScaRf1960MathOptInventory}.  It is optimal for single-echelon
problems with fixed ordering costs under mild conditions.  The
\textbf{$(R,Q)$ policy} places a fixed quantity $Q$ whenever the
inventory position reaches the reorder point $R$, which aligns naturally
with discrete vehicle load sizes.

\subsection{Economic Order Quantity (EOQ)}

The \textbf{EOQ model} \citep{Harris1913} minimises the sum of ordering
and holding costs in a deterministic, continuous-review setting:

\begin{equation}
  Q^* = \sqrt{\frac{2Kd}{h}}
\end{equation}

where $K$ is the fixed cost per order, $d$ is the demand rate, and $h$
is the holding cost rate.  EOQ is a useful reference for calibrating
order quantities in the $(R,Q)$ policy but ignores service levels and
capacity constraints.

%% -----------------------------------------------------------------------
\section{Multi-Echelon Optimality --- Clark and Scarf (1960)}
\label{sec:lit:clarkscarf}

\citet{ClarkScarf1960} established the foundation of MEIO by proving that,
for a serial network with backlogging, independent demands, and infinite
supplier capacity, the \textbf{echelon base-stock policy} is optimal.

The key insight is the \emph{echelon inventory position} (EIP): rather
than conditioning the replenishment decision at echelon $i$ on only the
local on-hand and in-transit stock, the policy conditions on the stock
at echelon $i$ \emph{plus} all downstream echelons.  This accounts for
the fact that inventory at echelon $i$ will eventually flow downstream
and must satisfy downstream demand as well as local demand.

The Clark-Scarf decomposition reduces the multi-echelon problem to a
sequence of single-echelon newsvendor problems, one per echelon, solved
in downstream-to-upstream order.  The solution is optimal under the
following assumptions:
\begin{itemize}
  \item Serial (chain) topology — no divergent DCs serving multiple retailers.
  \item Stationary, independent demands at each echelon.
  \item No capacity constraints on orders or vehicles.
  \item Backlogging (no lost sales).
\end{itemize}

When these assumptions are violated — as they are in the divergent 1N3
topology with vehicle capacity and real M5 demand — the Clark-Scarf
solution is no longer guaranteed to be optimal.  This gap motivates
the simulation-optimisation approach taken in this study.

%% -----------------------------------------------------------------------
\section{Periodic-Review and Cyclic Policies}
\label{sec:lit:periodic}
\vspace{1em}

\subsection{$(R,S)$ periodic-review policy}

The \textbf{$(R,S)$ policy} reviews inventory every $R$ periods and
orders up to level $S$.  It reduces the administrative overhead of
continuous-review policies and is widely used in practice.  Optimal
values of $R$ and $S$ can be found by exhaustive search or approximation
\citep{Silver1978}.

\subsection{$(k,m)$-cycle policies}

\citet{Roundy1985} showed that \textbf{$(k,m)$-cycle policies} — which
allow at most $k$ orders within an $m$-period cycle — can achieve
near-optimal cost in multi-item inventory systems while enabling
\emph{co-ordinated replenishment}.  By fixing the days on which each
item is reviewed, transport planners can schedule full-truckload
deliveries in advance, reducing the empty-vehicle problem that plagues
non-coordinated ordering.

%% -----------------------------------------------------------------------
\section{The Bullwhip Effect}
\label{sec:lit:bullwhip}

\citet{LeeEtAl1997} identified and formalised the \textbf{bullwhip
effect}: the tendency for demand variability to amplify as orders
propagate upstream in a supply chain.  They attributed bullwhip to four
root causes: demand signal processing (over-reaction to demand
fluctuations), order batching, price fluctuations, and rationing games.

The most widely used empirical metric for bullwhip is the ratio of order
variability to demand variability \citep{ChenEtAl2000}, applied
\emph{per echelon}:

\begin{equation}
  \text{BW}_{n,s} \;=\;
  \frac{\text{CV}^2(\text{orders placed by node } n \text{ for SKU } s)}
       {\text{CV}^2(\text{demand faced by node } n \text{ for SKU } s)}
\end{equation}

A ratio greater than one at an upstream node indicates amplification.
A subtlety, important for multi-echelon networks, is the choice of
denominator: the ``demand faced by node $n$'' is the \emph{external}
customer demand only when $n$ is a retailer.  At an internal node such
as a warehouse, the relevant denominator is the time series of
\emph{aggregated child orders} arriving at $n$, not the raw external
demand at any one downstream retailer \citep{LeeEtAl1997}.  Failing to
aggregate produces an inflated, mis-attributed ratio at upstream nodes.
This metric is computed per (node, SKU) on the post-warm-up window of
each simulation run (Section~\ref{sec:meth:bullwhip}).

%% -----------------------------------------------------------------------
\section{Transport Consolidation in Inventory Models}
\label{sec:lit:transport}

Classical MEIO models treat transport cost as a linear function of
quantity shipped, ignoring the step-function cost structure of
vehicle-based transport.  In practice, a truck that is 50\% full
and a truck that is 100\% full cost nearly the same to operate; the
cost-per-unit is therefore roughly twice as high for the half-full truck.

Research on coordinated inventory and transport (see, e.g.,
\citealt{JagermanEtAl1990} and related work) showed that the
interaction between inventory and transport decisions can be significant:
coordinating order quantities with vehicle load sizes allows suppliers
to dispatch fuller trucks, substantially reducing per-unit transport
cost in realistic settings.

Despite this, the dominant MEIO literature --- including Clark-Scarf,
Graves-Willems, and most textbook treatments --- ignores vehicle capacity
and consolidation.  This creates the gap that the joint optimisation
framework in this study addresses: when transport cost dominates the budget,
\emph{inventory and transport decisions must be optimised jointly}.

%% -----------------------------------------------------------------------
\section{Simulation-Optimisation Methods}
\label{sec:lit:simopt}

When the performance surface of a stochastic or complex system cannot be
expressed analytically, \textbf{simulation-optimisation} uses a simulator
as an oracle and an optimisation algorithm to search the parameter space
\citep{FuEtAl2005}. Various optimisation approaches are discussed here.

\subsection{Classical gradient-free approaches}

Genetic algorithms (GAs) \citep{Goldberg1989} have been applied to MEIO
parameter tuning since the 1990s.  They are flexible and parallelisable
but typically require thousands of function evaluations to converge and
are not sample-efficient for smooth, low-dimensional surfaces.
\textbf{Bayesian optimisation} \citep{ShahriariEtAl2016} addresses the
sample-efficiency gap by fitting a Gaussian-process surrogate and using an
acquisition function (e.g.\ expected improvement) to select the next
evaluation point.  It is highly sample-efficient for smooth functions in up
to $\sim\!20$ dimensions but scales cubically in the number of observations,
making it slow for the 30-dimensional joint space studied here.

\subsection{Tree-structured Parzen Estimator (TPE)}

The \textbf{TPE algorithm} \citep{BergstraEtAl2011} replaces the
Gaussian-process surrogate with two separate kernel density estimators:
one for the distribution of parameters that led to good results ($\ell(x)$)
and one for the rest ($g(x)$).  The acquisition function is
$\ell(x) / g(x)$, maximised by drawing many samples from $\ell(x)$ and
selecting the one with the highest ratio.

TPE is well-suited to this study's optimisation problem because:
\begin{itemize}
  \item It handles mixed-integer variables (integer base-stock levels and
        continuous dispatch thresholds) natively.
  \item It scales approximately linearly in dimensionality, remaining
        tractable for the 30-dimensional joint search space (25 stock
        levels + 5 dispatch thresholds).
  \item It is implemented in \textsc{Optuna} \citep{AkibaEtAl2019},
        an open-source framework with support for parallelism and
        constraint handling.
\end{itemize}

\subsection{Multi-objective optimisation: NSGA-II}

When the goal is not to satisfy a hard service-level constraint but to
\emph{trace the trade-off} between cost and service level, a
multi-objective evolutionary algorithm becomes appropriate.
\textbf{NSGA-II} \citep{DebEtAl2002} maintains a population of candidate
solutions and uses non-dominated sorting and crowding-distance
diversification to converge on the Pareto front in a single run.  In
this study, NSGA-II — implemented via Optuna's NSGAIISampler —
minimises $\bigl(C_{\text{total}},\;-\hat{\alpha}\bigr)$ jointly
(i.e.\ minimises total cost while maximising fill rate $\hat{\alpha}$)
and returns the non-dominated set as the Pareto front.  This produces a
smooth, monotone cost--service curve in a single run, replacing the
earlier constraint-sweep approach in which the joint optimiser was
re-run independently at each of several fill-rate targets; that
constraint-sweep is retained as a sanity-check overlay
(see Section~\ref{sec:opt:pareto}).

\subsection{Neural and reinforcement-learning approaches}
\label{sec:lit:drl}

Recent years have seen growing interest in applying \textbf{deep
reinforcement learning} (DRL) to inventory control.
\citet{HubbsEtAl2020} introduce \textsc{OR-Gym}, an
OpenAI-Gym-compatible benchmark suite that includes single-echelon and
serial multi-echelon environments; they show that policy-gradient methods
(PPO, A2C) can match or exceed newsvendor-optimal policies in these
settings.
\citet{GijsbrechtsEtAl2022} demonstrate that DRL agents can learn
near-optimal policies for the lost-sales problem with positive lead
times — a notoriously intractable case for exact methods.
\citet{PerezEtAl2021} extend DRL to stochastic serial systems, and
\citet{BouteEtAl2022} provide a practitioner-oriented survey of DRL for
supply-chain management.

Despite this progress, DRL approaches present several challenges for the
problem studied here.  First, the action space of the current framework
is \emph{mixed-integer}: base-stock levels are discrete (units) while
dispatch thresholds are continuous.  Policy-gradient methods require
differentiable action distributions, which are poorly suited to large
integer spaces without continuous relaxation.  Second, DRL requires
extensive environment interaction to learn a policy from scratch; the
30-dimensional joint search space in this study would require
substantially more than $10^{3}$ simulator evaluations to converge,
rendering training time prohibitive within the scope of this study.  Third, the
parameters optimised by this framework are physically interpretable:
base-stock levels map directly to safety-stock decisions and dispatch
thresholds to consolidation strategy, supporting managerial interpretation
of the result.  TPE was therefore chosen over DRL for its
sample-efficiency, mixed-integer native support, and the interpretability
of its output.

%% -----------------------------------------------------------------------
\section{Existing MEIO Software}
\label{sec:lit:software}

Several simulation and optimisation tools are available for supply-chain
modelling, but none satisfies all requirements of this study simultaneously.

\begin{description}
  \item[\textsc{SimPy}] A Python discrete-event simulation library.
        General-purpose and well-maintained, but provides no MEIO domain
        primitives — no nodes, policies, or transport layer.  Building an
        MEIO simulator from SimPy requires implementing all domain logic
        from scratch.

  \item[\textsc{AnyLogic}] A commercial simulation platform with a
        supply-chain library.  Powerful and well-documented, but
        proprietary, expensive, and not scriptable for large-scale
        parameter sweeps.

  \item[\textsc{simopt}] An open-source Python library for
        simulation-optimisation \citep{SimOptLib2022}.  Provides a
        collection of problem instances and solvers, but does not
        include MEIO-specific features such as multi-SKU demand,
        vehicle capacity, or transport consolidation.

  \item[\textsc{Optuna}] An open-source Bayesian/TPE optimisation
        framework \citep{AkibaEtAl2019}.  Used as the outer
        optimisation loop in this study; does not include simulation
        functionality.
\end{description}
The tools surveyed above cover different parts of the problem well, but none combines
an MEIO-specific simulation engine, a joint inventory-and-transport optimiser, and
support for arbitrary DAG topologies and vehicle-capacity constraints in a single,
open, scriptable package.  This work builds such a framework and uses it to study
the empirical value of joint optimisation on realistic demand data.
